\section{Limitations}
\label{sec:limitations}

We are deliberate about what this preliminary draft does and does not establish.

\begin{itemize}[leftmargin=1.4em,itemsep=0.2em]
  \item \textbf{Small evaluation set.} All quantitative results are on a single
        held-out set of $269$ trajectories ($130$ attacks / $139$ benign).
        Confidence intervals around PR@FSR are wide at this scale, and the
        headline operating point can move by a few traces. Larger and more
        diverse evaluation sets are needed before treating these numbers as
        benchmark figures.
  \item \textbf{Single training run.} We report one v0.19 run. We do not yet
        have seed variance, and we have not run the ablations that would isolate
        the contribution of each design choice (dense per-step supervision,
        $W_{\textsc{STOP}}$ up-weighting, anomaly truncation, left-truncation).
        \TODO{run ablations and report mean$\pm$std over seeds.}
  \item \textbf{ATBench false-stop floor.} The residual false-stop rate is
        concentrated almost entirely in one benign source (ATBench,
        FSR $\approx 0.065$ vs. $\approx 0.000$ elsewhere). This caps how tight
        the operating point can be set and is the current binding constraint;
        it is a data-coverage problem on that source rather than a method
        problem, but until it is fixed the practical FSR floor is set by it.
  \item \textbf{Per-family imbalance.} Detection is near-perfect on shaped
        families that dominate training but weak on under-represented or
        stylistically distinct families (ATBench-native, MCP, DNS exfiltration).
        Reported aggregate prevention therefore over-weights the well-covered
        families.
  \item \textbf{No latency numbers yet.} The ``cheap enough to gate every
        step'' claim is architectural (sub-2B, single forward pass on the common
        path), but we have not yet measured per-check latency on GPU or CPU.
        \TODO{measure p50/p99 latency per check.}
  \item \textbf{No quantitative cross-backend transfer.} We argue
        canonicalization makes one model framework-agnostic and observe transfer
        qualitatively, but we have not measured the train-on-A / test-on-B
        prevention drop. \TODO{measure cross-backend transfer.}
  \item \textbf{Threshold calibration is per-deployment.} The operating
        threshold $\tau\approx 0.92$ was selected on held-out data; a new
        deployment with a different benign distribution would need to
        recalibrate, and we have not characterized how transferable $\tau$ is.
  \item \textbf{Evasion.} As a learned distributional monitor, IAM is in
        principle evadable by attacks crafted to look in-distribution. We have
        not yet run an adaptive-attacker evaluation. \TODO{adaptive evasion
        study.}
  \item \textbf{Scope.} IAM complements, and does not replace, model
        alignment (Section~\ref{sec:intro}). For strongly-aligned backends its
        marginal value is small; its intended regime is weakly-aligned
        deployments.
\end{itemize}
